%! The Language of Variation — Unit 1, Lesson 1.1 — Group Activity (Answer Key)
% GRADUAL RELEASE: the "you do together" block — 16 minutes in groups of 3-4, followed by an
% 8-minute whole-class debrief that the Debrief box at the end captures.
% ONE activity for the entire class — there are no tiers and no differentiated versions.
% The vocabulary was GIVEN in the guided notes, so it is used freely here. The work is applying
% the amount test to a register students have not seen, and defending each call.
% CRUX: scenario 2. Four columns written entirely in digits, only two of which are quantitative.
% The average of two permit numbers computes cleanly and means nothing.
%
% PAGE LOCKSTEP with activity_key/: \answerspace{H}{} reserves exactly H in the blank and prints
% the red answer in the key. Keep the macro, every height, and every \boxguard byte-identical.
\documentclass[10pt]{article}
\usepackage{apstats-article}
\usepackage{apstats-key}   % pulls in -boxes; do NOT also load -boxes
\newcommand{\answerspace}[2]{\par\nopagebreak\noindent\begin{minipage}[t][#1][t]{\linewidth}%
  \color{keyred}\bfseries #2\end{minipage}\par}

\begin{document}

\pageheader{Unit 1, Lesson 1.1}{Group Activity: The Trailhead Register --- Answer Key}
% NAMESTRIP: no \namedateperiod here — mirrors the blank.

\vspace{0.07in}

\boxguard[10]
\begin{headlinebox}{sky}
\small
\textbf{The Trailhead Register.} Every party that hikes at Ross Hollow State Park signs the
register at the trailhead and signs out again on the way home. Work through both problems with
your group, using today's vocabulary. \textbf{Every classification needs a reason} --- if
somebody in the group cannot say \emph{why} a column is C or Q, nobody writes it down yet.
\end{headlinebox}

\vspace{0.07in}

\boxguard[28]
\begin{scenariobox}[1. Sorting the Register]{navy}
\small
\renewcommand{\arraystretch}{1.1}
\begin{tabularx}{\linewidth}{@{} l Y Y Y Y Y @{}}
\rowcolor{sky}
\textbf{Party} & \textbf{Trail} & \textbf{Group size} & \textbf{Time out (min)} &
\textbf{Dog along} & \textbf{Weather} \\
\hline
Alvarez  & Ridge Loop   & 4 &  95 & Yes & Sunny  \\
Boone    & Cedar Falls  & 2 & 140 & No  & Cloudy \\
Chen     & Ridge Loop   & 6 &  80 & Yes & Sunny  \\
Delgado  & Summit Trail & 3 & 215 & No  & Rain   \\
Egan     & Cedar Falls  & 2 & 155 & Yes & Cloudy \\
\hline
\end{tabularx}

\vspace{6pt}
\begin{enumerate}[label=\textbf{\alph*.}, leftmargin=1.8em, itemsep=5pt]
  \item One \textbf{row} of this register describes one \ans{individual --- one party}, and the register
        records \ans{5} variables about each.
  \item Classify each recorded variable. Write \textbf{C} or \textbf{Q}.

        \vspace{3pt}
        \renewcommand{\arraystretch}{1.3}
        \begin{tabularx}{\linewidth}{@{} Y Y Y Y Y @{}}
        \textbf{Trail} & \textbf{Group size} & \textbf{Time out} & \textbf{Dog along} &
        \textbf{Weather} \\
        \ans{C} & \ans{Q} & \ans{Q} & \ans{C} & \ans{C} \\
        \end{tabularx}
  \item Pick \textbf{one} column you marked \textbf{Q} and \textbf{one} you marked \textbf{C}.
        For each, run the amount test out loud in your group, then write the reason here.
        \answerspace{1.5cm}{Answers vary. \emph{Time out} is Q --- 215 minutes is an amount, and
        215 really is 120 more than 95. \emph{Trail} is C --- ``Ridge Loop'' names which trail,
        and Summit Trail is not ``more'' than Cedar Falls.}
  \item Adding up a column and dividing by 5 gives the ranger something useful for exactly the
        columns you marked \textbf{Q}. Write down the nonsense you would get if you tried it on
        \textbf{Weather}, then say what the ranger should do with that column instead.
        \answerspace{1.5cm}{Sunny $+$ Cloudy $+$ Sunny $+$ Rain $+$ Cloudy has no total, so
        there is no ``average weather'' to divide out. The ranger should \textbf{count} instead:
        2 sunny days, 2 cloudy, 1 rain --- or report them as percents of the five parties.}
\end{enumerate}
\end{scenariobox}

\vspace{0.07in}

\boxguard[26]
\begin{scenariobox}[2. Four Columns Written in Digits]{navy}
\small
Next season the ranger adds four more columns. Here are the first three rows.

\vspace{5pt}
\renewcommand{\arraystretch}{1.1}
\begin{tabularx}{\linewidth}{@{} l Y Y Y Y @{}}
\rowcolor{sky}
\textbf{Party} & \textbf{Permit \#} & \textbf{Parking spot} & \textbf{Home ZIP} &
\textbf{Water bottles} \\
\hline
Alvarez & 2041 & 17 & 80211 & 6 \\
Boone   & 6318 & 04 & 80304 & 2 \\
Delgado & 1150 & 11 & 80211 & 3 \\
\hline
\end{tabularx}

\vspace{6pt}
\begin{enumerate}[label=\textbf{\alph*.}, leftmargin=1.8em, itemsep=5pt]
  \item Every one of these four columns is written entirely in digits. Classify each anyway.

        \vspace{3pt}
        \renewcommand{\arraystretch}{1.3}
        \begin{tabularx}{\linewidth}{@{} Y Y Y Y @{}}
        \textbf{Permit \#} & \textbf{Parking spot} & \textbf{Home ZIP} & \textbf{Water bottles} \\
        \ans{C} & \ans{C} & \ans{C} & \ans{Q} \\
        \end{tabularx}
  \item \textbf{Two of your four calls are the ones people get wrong.} Name them, and explain
        in one sentence each what those digits are doing instead of measuring.
        \answerspace{1.8cm}{\emph{Permit \#} and \emph{Home ZIP} (and \emph{Parking spot}).
        A permit number identifies which party pulled it --- the average, 4179.5, is nobody's
        permit. A ZIP code names a postal area; 80211 is not ``more'' than 80304.}
  \item The season report says \textbf{``62\% of parties hiked with a dog.''} That number came
        out of the \emph{Dog along} column. Does producing a percent make \emph{Dog along}
        quantitative? Explain.
        \answerspace{1.8cm}{No. Counting a categorical variable \emph{produces} numbers --- that
        is how categorical variables are summarized. The 62\% is a numerical summary; the values
        in the column are still Yes and No.}
  \item Next season the ranger stops recording \emph{Time out} in minutes and records
        \emph{Short}, \emph{Medium}, or \emph{Long} instead. Name the kind of variable it
        becomes, and one thing the ranger can no longer compute.
        \answerspace{1.6cm}{It becomes \textbf{categorical} --- the values are now group labels.
        The ranger can no longer compute the mean time on the trail, because ``Medium'' does not
        carry the 140 minutes it came from. Grouping runs one way only.}
\end{enumerate}
\end{scenariobox}

\vspace{0.07in}

% ── DEBRIEF (8 min, whole class — filled together after the groups report out) ─
\boxguard[18]
\begin{reflectionbox}
\small
\textbf{Debrief --- we fill this in together.}
\begin{enumerate}[leftmargin=1.9em, itemsep=6pt]
  \item State the amount test in one sentence, without using the word ``digits.''
        \answerspace{1.0cm}{Ask whether the value records an amount --- how much or how many.
        If it only identifies or labels the individual, the variable is categorical.}
  \item Give one variable that is quantitative and one that is categorical \emph{that were not
        on this register}, and say how you know.
        \answerspace{1.2cm}{Answers vary: \emph{height in inches} is Q --- 68 inches is an amount.
        \emph{Eye color} is C --- brown names a category and has no total.}
  \item Why can grouping never be undone?
        \answerspace{1.0cm}{Because the label keeps only which bin the value fell into, not the
        value. ``Long'' cannot be turned back into 215 minutes.}
\end{enumerate}
\end{reflectionbox}

\end{document}
